\section{Introduction}
\label{intro:section}

An entangled projective measurement is constrained both by the orthogonality of its global basis states and by the geometry of their local reductions. The elegant joint measurement (EJM) brings these two features together: the local Bloch vectors occupy the vertices of a regular tetrahedron, while the corresponding global states form an orthonormal basis. Introduced in the study of quantum networks~\cite{Gisin2017}, the EJM at $n=2$ belongs to a continuous family connecting it to a Bell-state measurement~\cite{Tavakoli2021}. Varying the tetrahedral radius changes the single-qubit reduced spectrum and hence the entanglement across each single-qubit-versus-rest cut. Extending this control to more qubits requires the local geometries of all measurement outcomes to remain compatible with global orthogonality.

Pauli-orbit constructions provide a useful setting for this problem~\cite{PauwelsPauli2026,PauwelsEJM2026}. In the associated phase normal form, every diagonal unitary generates a complete orthonormal basis; the remaining constraints concern the local Bloch vectors. Pauwels and Gisin gave an explicit reference EJM for every $n\ge2$, with common radius
\begin{equation}
 r_E(n)=\frac{\sqrt3}{2^{n-1}},
 \label{intro:reference-radius}
\end{equation}
and an even-qubit phase family with radius $r_E(n)\cos t$ on $|t|\le\pi/2$~\cite{PauwelsGisin2026}. Their construction already identifies interference involving a single computational-basis phase correction as the source of the local signal. Their odd-qubit analysis leaves an explicit tunable family through the specified reference point open for $n\ge5$, and identifies a locally isolated reference point at $n=3$ from a numerical Jacobian rank. Other symmetric multiqubit constructions permit broader local geometries~\cite{Ding2025,He2026}; here the target remains one fixed regular tetrahedron at each qubit, with a common radius across qubits and outcomes.

An earlier version of the present work supplied an explicit enhanced family for every odd $n\ge5$, with radius $r_E(n)(1-\sin t)$~\cite{Zhang2026Odd}. The present construction removes the odd--even distinction: a different Boolean kernel, depending on four binary features, gives the same radius law for every integer $n\ge5$. It recovers the reference measurement column by column at $t=0$ and continuously connects zero local polarization to twice the reference radius. This enhancement raises a natural question about its limit. We answer it for all two-valued real phase generators through the reference measurement by classifying the admissible radius laws. The doubled radius is optimal in this model, which includes generators beyond the four-feature construction.

The unified construction also makes it possible to compare the new enhancement with the known even-qubit cosine deformation within one phase representation. We examine whether both can be varied independently while preserving the local tetrahedra. For every even $n\ge8$, the answer is an explicit two-parameter family of complete measurements containing the two deformations as exact coordinate sections. At $n=6$, mixed interference obstructs this extension; in the corresponding four-color phase model, the local solutions consist of two crossing branches.

The lower-dimensional cases clarify why a straight phase path can cease to be sufficient. No two-valued phase line enhances the reference radius for $n=2,3,4$. At $n=3$, we strengthen the numerical evidence of Pauwels and Gisin to an exact local-rigidity certificate. At $n=4$, even allowing arbitrary real generators leaves every admissible straight line nonenhancing, but an explicit curved path increases the radius to $(1+1/\sqrt2)r_E(4)$. Thus the obstruction at $n=4$ can be overcome by changing the shape of the path within the same phase representation.

Section~II develops the phase framework that reduces these questions to local interference conditions. Sections~III and IV give the unified construction and its sharp Boolean-model bound. Section~V establishes its relation to the even-qubit family, and Sec.~VI explains the low-dimensional obstructions and the curved construction at $n=4$. The appendixes provide the full proofs.
